\section{Reproduction and repository integration}
\label{app:reproduction}

The archive is a self-contained research continuation, with the following
layout relative to its top-level directory:
\begin{center}
\begin{tabular}{p{0.28\textwidth}p{0.63\textwidth}}
\toprule
Path & Contents \\
\midrule
\path{article/} & This PDF, its main \TeX{} source, sections, bibliography,
generated tables, vector figures, and build instructions. \\
\path{code/fast/} & Runnable maintained source snapshot with the new modules,
tests, and dependency-free demonstration. \\
\path{code/reports/} & Minimal pinned archived oracle dependencies needed by
the complete maintained test suite. \\
\path{reference/} & Exact pinned historical interval producer used for
paired baseline measurements. \\
\path{experiments/} & Timing driver, independent component audit and frozen
trefoil fixture, and table/figure generator. \\
\path{results/} & Final raw measurements, finite-audit records, test log,
weighted-test summary, and concise demonstration output. \\
\path{integration.patch} & Changes to production code, tests, demonstration,
and the maintained README, with paths relative to the ProveIt repository. \\
\path{PROVENANCE.json} & Source pin, file identities, environment information,
and evidence-to-command mapping. \\
\bottomrule
\end{tabular}
\end{center}

The runtime source and demonstration require Python 3.10 or later; the
recorded environment uses 3.12.14.  The complete independent surface audit
requires Regina; the project already specifies its optional Regina range
in \path{code/fast/pyproject.toml}.  Plot regeneration requires Matplotlib.
PDF compilation uses a standard \LaTeX{} installation with pdf\LaTeX{},
Bib\TeX{}, Latin Modern, \code{microtype}, \code{cleveref}, and the packages
listed in the main source.  No external file download is required to read
the article, run the demonstration, or compare the supplied interval code.

From the archive root, the convenience driver exposes the independent jobs:
\begin{verbatim}
python reproduce.py demo
python reproduce.py tests
python reproduce.py audit
python reproduce.py benchmark
python reproduce.py figures
python reproduce.py article
\end{verbatim}
Regenerated outputs are written under \path{rerun/}, leaving the recorded
evidence intact.  Timings should be collected without concurrent heavy
jobs.  Repeated runs may change elapsed times while leaving operation
counts, component predicates, and certificate checks identical.

The benchmark command compares the delivered source to
\path{reference/interval_orbits.py}; it does not fetch a moving branch.
The independent normal audit reads the adjacent frozen trefoil fixture and
enumerates only the explicitly described finite families.  The weighted
test suite writes its small diagnostic JSON into the research folder, which
the convenience driver copies into the rerun results.  The driver reports
subprocess failure instead of silently continuing to a successful summary.

Before integrating code into another checkout, run
\begin{verbatim}
git apply --check /path/to/integration.patch
git apply /path/to/integration.patch
\end{verbatim}
from the ProveIt root.  The patch targets the pinned commit; later
concurrent changes may require ordinary review and conflict resolution.
The archive's source snapshot is useful for reproduction, but replacing
the whole current \path{fast/} directory with it would discard later work.
For the manuscript, copy \path{article/} as a complete subdirectory, for
example under
\path{Topology/UnknotRecognition/synthesis/component_certificates_20261009/}.
Its relative section, table, figure, and bibliography paths then remain valid.

The new production sources follow the maintained MIT No Attribution license.
Existing source and archived oracle license notices are retained.  The
bundle references third-party papers bibliographically; it does not
redistribute their downloaded PDFs.  The recorded raw results are evidence
for this continuation, while the proofs and source comments state the
generality and limits of each claim.
